\documentclass[11pt,a4paper,UTF8]{ctexart}
\usepackage{geometry}
\geometry{left=18mm,right=18mm,top=24mm,bottom=24mm,headheight=22pt,footskip=12mm}
\usepackage{amsmath,amssymb,mathtools,bm,esint,extarrows}
\usepackage{xcolor}
\usepackage{tcolorbox}
\tcbuselibrary{skins,breakable}
\usepackage{fancyhdr}
\usepackage{titlesec}
\usepackage{hyperref}
\usepackage{tikz}
\usepackage{booktabs}
\usepackage{tabularx}

% --- Color Palette (Purple & DeepPink Academic Theme) ---
\definecolor{DeepPurple}{HTML}{4C1D95}
\definecolor{TitlePurple}{HTML}{581C87}
\definecolor{SubPurple}{HTML}{6B21A8}
\definecolor{DeepPink}{HTML}{FF1493}
\definecolor{LightPinkBg}{HTML}{FFF5F8}
\definecolor{LightPurpleBg}{HTML}{F8F5FF}
\definecolor{GoldAccent}{HTML}{D97706}
\definecolor{DarkText}{HTML}{1F2937}

\hypersetup{
    colorlinks=true,
    linkcolor=TitlePurple,
    citecolor=DeepPink,
    urlcolor=SubPurple,
    pdfborder={0 0 0}
}

% --- Typography & Spacing ---
\linespread{1.3}
\setlength{\parskip}{0.35em plus 0.1em minus 0.1em}
\setlength{\parindent}{0pt}

% --- Section Titles Styling ---
\titleformat{\section}
  {\Large\bfseries\color{TitlePurple}}{\thesection}{1em}{}
  [\vspace{1mm}{\color{TitlePurple!30}\hrule height 1pt}]
\titleformat{\subsection}
  {\large\bfseries\color{SubPurple}}{\thesubsection}{0.8em}{}
\titleformat{\subsubsection}
  {\normalsize\bfseries\color{DeepPurple}}{\thesubsubsection}{0.6em}{}

% --- tcolorbox Environments ---
% 1. Theorem / Formula / Method Box (Deep Purple)
\newtcolorbox{thmbox}[1][]{
    enhanced,
    breakable,
    colback=LightPurpleBg,
    colframe=TitlePurple,
    arc=3pt,
    boxrule=0pt,
    leftrule=3.5pt,
    left=10pt,
    right=10pt,
    top=8pt,
    bottom=8pt,
    title=#1,
    coltitle=white,
    colbacktitle=TitlePurple,
    fonttitle=\bfseries\small,
    attach boxed title to top left={yshift=-2mm, xshift=4mm},
    boxed title style={boxrule=0pt, arc=2pt}
}

% 2. Solution / Formula / Derivation Box (DeepPink)
\newtcolorbox{solbox}[1][]{
    enhanced,
    breakable,
    colback=LightPinkBg,
    colframe=DeepPink,
    arc=3pt,
    boxrule=0pt,
    leftrule=3.5pt,
    left=10pt,
    right=10pt,
    top=8pt,
    bottom=8pt,
    title=#1,
    coltitle=white,
    colbacktitle=DeepPink,
    fonttitle=\bfseries\small,
    attach boxed title to top left={yshift=-2mm, xshift=4mm},
    boxed title style={boxrule=0pt, arc=2pt}
}

% 3. Learning Goals / Summary Box
\newtcolorbox{targetbox}[1][]{
    enhanced,
    breakable,
    colback=LightPurpleBg,
    colframe=DeepPurple,
    arc=3pt,
    boxrule=1pt,
    left=10pt,
    right=10pt,
    top=8pt,
    bottom=8pt,
    title=#1,
    coltitle=white,
    colbacktitle=DeepPurple,
    fonttitle=\bfseries\small,
    attach boxed title to top left={yshift=-2mm, xshift=4mm},
    boxed title style={boxrule=0pt, arc=2pt}
}

\begin{document}

% ------------------- 讲义封面 (Cover Page) -------------------
\begin{titlepage}
\thispagestyle{empty}
\begin{tikzpicture}[remember picture,overlay]
    \fill[DeepPurple] (current page.north west) rectangle ([yshift=-7cm]current page.north east);
    \draw[DeepPink, line width=3.5pt] ([yshift=-7cm]current page.north west) -- ([yshift=-7cm]current page.north east);
    \draw[GoldAccent, line width=1pt] ([yshift=-7.12cm]current page.north west) -- ([yshift=-7.12cm]current page.north east);
    
    \node[anchor=north east, opacity=0.12, text=white] at ([xshift=-1.5cm, yshift=-0.5cm]current page.north east) {
        \fontsize{70}{70}\selectfont\bfseries 2027
    };
\end{tikzpicture}

\vspace*{0.8cm}
\begin{center}
    {\color{white}\fontsize{26}{32}\selectfont\bfseries 2027 考研数学(二) 核心讲义}\\[0.6cm]
    {\color{white}\fontsize{13}{18}\selectfont\bfseries 全国硕士研究生招生考试 · 统考数学(二) 核心精讲全书}\\[1.5cm]
\end{center}

\vspace{3.5cm}
\begin{center}
    {\color{GoldAccent}\large\bfseries $\blacktriangleright$ 基础强化 · 核心题解 · 考点深水区 $\blacktriangleleft$}\\[0.8cm]
    {\color{TitlePurple}\fontsize{24}{28}\selectfont\bfseries 第十章 一元积分学的应用(几何)}\\[0.6cm]
    {\color{DeepPink}\large\bfseries —— 体系化理论精讲与公式全景 ——}\\[1.5cm]
\end{center}

\vfill

\begin{center}
\begin{tcolorbox}[
    colback=white,
    colframe=DeepPurple,
    width=0.88\textwidth,
    arc=4pt,
    boxrule=1.2pt,
    halign=center,
    drop shadow=gray!30
]
    \vspace{3mm}
    {\bfseries\large 编著团队：EscoffierZhou}\\[2.5mm]
    {\bfseries\color{TitlePurple} 目标院校：中国科学院大学 (UCAS) 计算机专硕 (22408)}\\[2.5mm]
    {\small\color{gray} 备考铁律：手算真题数字 · 错题白纸重做 · 408 客观题为王}\\[2.5mm]
    {\small 编制版本：2027 届首发版 · \today}
    \vspace{2mm}
\end{tcolorbox}
\end{center}
\vspace{0.8cm}
\end{titlepage}

% ------------------- 目录与页面主体配置 -------------------
\pagestyle{fancy}
\fancyhf{}
\fancyhead[L]{\small\color{gray}2027 考研数学(二) · 核心讲义系列}
\fancyhead[R]{\small\color{TitlePurple}\nouppercase{\leftmark}}
\fancyfoot[C]{\small\color{gray}— 第 \thepage\ 页 —}
\fancyfoot[R]{\small\color{gray}UCAS 22408}
\renewcommand{\headrulewidth}{0.5pt}

\pagenumbering{Roman}
\tableofcontents
\newpage
\pagenumbering{arabic}


\begin{targetbox}[本章复习目标与核心知识图谱]
\begin{itemize}
  \item 目的[1]:深刻理解微元分析法（“度量 vs 本质”与 $o(\Delta x)$ 高阶无穷小忽略准则），熟练掌握平面直角坐标、极坐标、参数方程三种体系下的平面图形面积微元构建与计算
\end{itemize}
\end{targetbox}

\begin{targetbox}[本章复习目标与核心知识图谱]
\begin{itemize}
  \item 目的[2]:彻底掌握旋转体体积两大经典微元（切片圆盘法/圆环法 Disc/Washer Method 与柱壳法 Cylindrical Shells Method），做到根据图形特征自由秒切
\end{itemize}
\end{targetbox}

\begin{targetbox}[本章复习目标与核心知识图谱]
\begin{itemize}
  \item 目的[3]:深刻掌握平面曲线弧长微元 $\mathrm{d}s$ 在直角坐标、参数方程、极坐标三种体系下的统一定义与精准计算
\end{itemize}
\end{targetbox}

\begin{targetbox}[本章复习目标与核心知识图谱]
\begin{itemize}
  \item 目的[4]:掌握旋转曲面的侧面积微元 $\mathrm{d}S = 2\pi r\,\mathrm{d}s$ 的本质机理与易错陷阱（绝对不可用 $\mathrm{d}x$ 替代弧长微元 $\mathrm{d}s$）
\end{itemize}
\end{targetbox}

\begin{targetbox}[本章复习目标与核心知识图谱]
\begin{itemize}
  \item 目的[5]:全面掌握由已知平行截面面积求立体体积的截面积分法，以及函数平均值定理及其变限微分方程逆问题
\end{itemize}
\end{targetbox}

\begin{targetbox}[本章复习目标与核心知识图谱]
\begin{itemize}
  \item 目的[6]:系统掌握平面均质薄板形心坐标公式、平面曲线质心公式及其与二重积分重心的微积分同构关系
\end{itemize}
\end{targetbox}


\section{1.微元分析法核心思想与平面图形的面积}

\begin{thmbox}[方法[1]: 微元法（元素法）的本质与“度量 vs 本质”辨析]

在高等数学与考研数学中，定积分的最核心应用在于“微元法”（又称元素法）。微元法是将所求几何量或物理量 $U$ 表达为定积分的有效工具。

\textbf{[1]} 微元法的标准实施三部曲：
\textbf{(1)}  \textit{}分割与近似（求微元 $\mathrm{d}U$）\textit{}：

在自变量变化区间 $[a, b]$ 上任取一个微小区间 $[x, x + \mathrm{d}x]$。设在这一小区间上，量 $U$ 的微小增量为 $\Delta U$。若能找到一个形如 $f(x)\,\mathrm{d}x$ 的连续微分表达式，使得：


\begin{equation*}
\Delta U = f(x)\,\mathrm{d}x + o(\mathrm{d}x)
\end{equation*}


其中 $o(\mathrm{d}x)$ 是当 $\mathrm{d}x \to 0$ 时比 $\mathrm{d}x$ 更高阶的无穷小量，则称 $\mathrm{d}U = f(x)\,\mathrm{d}x$ 为所求量 $U$ 的\textbf{微元}（或微分元素）。

\textbf{(2)}  \textbf{求和与取极限（积分离散化为整体）}：

对区间 $[a, b]$ 上的全部微元进行连续累加，即取定积分：


\begin{equation*}
U = \int_a^b \mathrm{d}U = \int_a^b f(x)\,\mathrm{d}x
\end{equation*}


\textbf{[2]} “度量 vs 本质”的哲学辨析与高阶无穷小舍弃原则：
\textbf{(1)}  \textbf{曲边梯形面积中的高阶误差}：

以曲边梯形为例，底边为 $\Delta x$，高为曲线高度。若用矩形面积 $f(x)\Delta x$ 近似梯形面积 $\Delta U$，误差为一个小三角形或小弯曲区域。由于曲线连续可导，该弯折区域的面积不超过 $\frac{1}{2}|\Delta f|\Delta x = \frac{1}{2}|f'(\xi)|(\Delta x)^2 = o(\Delta x)$。因此在极限求和 $\sum \Delta U$ 中，全部误差之和 $\sum o(\Delta x) \to 0$。用矩形微元 $f(x)\,\mathrm{d}x$ 完全精确。

\textbf{(2)}  \textbf{弧长与侧面积中的致命陷阱}：

在计算曲线弧长或旋转曲面侧面积时，若误以为倾斜曲线微元可以近似为水平线段 $\Delta x$，其误差为：


\begin{equation*}
\Delta s - \Delta x = \sqrt{(\Delta x)^2 + (\Delta y)^2} - \Delta x = \Delta x \left(\sqrt{1 + \left(\frac{\Delta y}{\Delta x}\right)^2} - 1\right)
\end{equation*}


此时当 $y' \neq 0$ 时，括号内为非零常数，误差是与 $\Delta x$ 同阶的\textbf{同阶无穷小}，而绝非高阶无穷小 $o(\Delta x)$！若直接忽略该误差，求和结果将彻底错误。因此，弧长必须采用斜边弦长 $\mathrm{d}s = \sqrt{\mathrm{d}x^2 + \mathrm{d}y^2}$ 作为微元。

\end{thmbox}
\begin{thmbox}[公式[1]: 直角坐标系下的平面图形面积]

设平面图形由连续曲线所围成，按投影方向划分为 $X$ 型与 $Y$ 型区域：

\textbf{[1]} $X$ 型区域（垂直条微元法）：
若平面区域 $D$ 由上边界 $y = y_2(x)$、下边界 $y = y_1(x)$ 以及两条垂直直线 $x = a, x = b$ ($a < b$) 围成，且 $y_2(x) \geqslant y_1(x)$：

取宽为 $\mathrm{d}x$ 的垂直狭长矩形微元，面积微元为：


\begin{equation*}
\mathrm{d}S = [y_2(x) - y_1(x)]\,\mathrm{d}x
\end{equation*}


平面图形总面积为：


\begin{equation*}
S = \int_a^b [y_2(x) - y_1(x)]\,\mathrm{d}x
\end{equation*}


\textbf{[2]} $Y$ 型区域（水平条微元法）：
若平面区域 $D$ 由右边界 $x = x_2(y)$、左边界 $x = x_1(y)$ 以及两条水平直线 $y = c, y = d$ ($c < d$) 围成，且 $x_2(y) \geqslant x_1(y)$：

取高为 $\mathrm{d}y$ 的水平狭长矩形微元，面积微元为：


\begin{equation*}
\mathrm{d}S = [x_2(y) - x_1(y)]\,\mathrm{d}y
\end{equation*}


平面图形总面积为：


\begin{equation*}
S = \int_c^d [x_2(y) - x_1(y)]\,\mathrm{d}y
\end{equation*}


\end{thmbox}

\begin{align*}
\begin{cases} x = x(t) \\ \\ y = y(t) \end{cases} \quad (\alpha \leqslant t \leqslant \beta)
\end{align*}
\begin{thmbox}[公式[2]: 参数方程下的平面图形面积]

若平面曲线由参数方程给出：


\textbf{[1]} 常规积分公式：
若参数 $t$ 从 $\alpha$ 增加到 $\beta$ 时，$x(t)$ 单调增加，且曲线在 $x$ 轴上方，面积微元为：


\begin{equation*}
\mathrm{d}S = y(x)\,\mathrm{d}x = y(t)x'(t)\,\mathrm{d}t
\end{equation*}


故平面图形面积为：


\begin{equation*}
S = \int_\alpha^\beta y(t)x'(t)\,\mathrm{d}t
\end{equation*}


若随着 $t$ 增大 $x(t)$ 递减，则加负号或取绝对值 $|x'(t)|$ 保持非负。

\textbf{[2]} 闭合曲线面积的对称式（格林公式导出）：
对于无自交的正向闭合光滑曲线 $L$（逆时针方向），围成区域 $D$ 的面积可表示为第二类曲线积分，根据格林公式：


\begin{equation*}
S = \iint_D \mathrm{d}x\mathrm{d}y = \frac{1}{2}\oint_L (x\,\mathrm{d}y - y\,\mathrm{d}x)
\end{equation*}


代入参数方程得高度对称的形式：


\begin{equation*}
S = \frac{1}{2}\int_\alpha^\beta [x(t)y'(t) - y(t)x'(t)]\,\mathrm{d}t
\end{equation*}


该公式在处理摆线拱、星形线、椭圆等参数曲线时极为优雅且不易混淆符号。

\end{thmbox}
\begin{thmbox}[公式[3]: 极坐标系下的平面图形面积]

在极坐标系 $(r, \theta)$ 下，平面图形由射线及极径曲线界定：

\textbf{[1]} 单曲线扇形微元：
由曲线 $r = r(\theta)$ 以及两条极角射线 $\theta = \alpha, \theta = \beta$ ($\alpha < \beta$) 所围成的曲边扇形：

取夹角为 $\mathrm{d}\theta$ 的微小扇形，由于半径为 $r(\theta)$，弧长为 $r(\theta)\,\mathrm{d}\theta$，小扇形面积近似为以 $r$ 为底、弧长为高的三角形面积：


\begin{equation*}
\mathrm{d}S = \frac{1}{2}r^2(\theta)\,\mathrm{d}\theta
\end{equation*}


图形总面积为：


\begin{equation*}
S = \frac{1}{2}\int_\alpha^\beta r^2(\theta)\,\mathrm{d}\theta
\end{equation*}


\textbf{[2]} 双曲线同向夹层面积：
若区域介于外边界 $r_2(\theta)$ 与内边界 $r_1(\theta)$ 之间 ($r_2(\theta) \geqslant r_1(\theta) \geqslant 0$)：


\begin{equation*}
\mathrm{d}S = \frac{1}{2}[r_2^2(\theta) - r_1^2(\theta)]\,\mathrm{d}\theta \implies S = \frac{1}{2}\int_\alpha^\beta [r_2^2(\theta) - r_1^2(\theta)]\,\mathrm{d}\theta
\end{equation*}


\end{thmbox}

\section{2.旋转体体积与已知平行截面面积的立体体积}

\begin{thmbox}[公式[4]: 绕坐标轴旋转的旋转体体积]

平面区域绕定直线旋转一周所形成的几何体称为旋转体。考研中根据切片方向分为“圆盘/圆环法”与“圆柱壳法”：

\textbf{[1]} 圆盘法与圆环法 (Disk / Washer Method，垂直切片)：
\textbf{(1)}  \textit{}绕 $x$ 轴旋转\textit{}：

由连续曲线 $y = f(x)$、$x$ 轴及垂线 $x = a, x = b$ ($a < b$) 围成的曲边梯形绕 $x$ 轴旋转：

垂直于 $x$ 轴切出的薄片为底半径为 $|y(x)|$、厚度为 $\mathrm{d}x$ 的圆盘，体积微元为：


\begin{equation*}
\mathrm{d}V_x = \pi y^2(x)\,\mathrm{d}x
\end{equation*}


总体积为：


\begin{equation*}
V_x = \pi \int_a^b y^2(x)\,\mathrm{d}x
\end{equation*}


若区域由两曲线 $y_2(x) \geqslant y_1(x) \geqslant 0$ 界定，则切片为圆环（Washer），体积为：


\begin{equation*}
V_x = \pi \int_a^b [y_2^2(x) - y_1^2(x)]\,\mathrm{d}x
\end{equation*}


\textbf{(2)}  \textit{}绕 $y$ 轴旋转（$Y$ 型截面）\textit{}：

若区域由 $x = g(y)$ 及水平线 $y = c, y = d$ 界定，绕 $y$ 轴旋转时垂直于 $y$ 轴切片：


\begin{equation*}
V_y = \pi \int_c^d x^2(y)\,\mathrm{d}y
\end{equation*}


\textbf{[2]} 圆柱壳法 (Cylindrical Shells Method，平行切片/套筒法)：
\textbf{(1)}  \textit{}绕 $y$ 轴旋转（自变量为 $x$）\textit{}：

曲边梯形 $0 \leqslant y \leqslant f(x), a \leqslant x \leqslant b$ ($0 \leqslant a < b$) 绕 $y$ 轴旋转：

在 $x$ 处取厚度为 $\mathrm{d}x$ 的竖向狭长条，绕 $y$ 轴旋转一周展开后为一个薄壁圆柱薄壳：

- 薄壳旋转半径为 $x$；
- 薄壳高度为 $f(x)$；
- 薄壳周长为 $2\pi x$；
- 薄壳壁厚为 $\mathrm{d}x$。

体积微元为薄壳侧面积乘壁厚：


\begin{equation*}
\mathrm{d}V_y = 2\pi x f(x)\,\mathrm{d}x
\end{equation*}


总体积为：


\begin{equation*}
V_y = 2\pi \int_a^b x f(x)\,\mathrm{d}x
\end{equation*}


\textbf{(2)}  \textit{}绕 $x$ 轴旋转（自变量为 $y$）\textit{}：

同理，由 $x = g(y), c \leqslant y \leqslant d$ ($0 \leqslant c < d$) 绕 $x$ 轴旋转：


\begin{equation*}
V_x = 2\pi \int_c^d y g(y)\,\mathrm{d}y
\end{equation*}


\textbf{[3]} 绕平行于坐标轴的直线旋转：
若旋转轴为水平线 $y = y_0$ 或垂直线 $x = x_0$：

- 圆盘法只需将回转半径修正为点到轴的垂直距离，如绕 $y = y_0$ 旋转：$\mathrm{d}V = \pi |y(x) - y_0|^2\,\mathrm{d}x$；
- 圆柱壳法只需将圆柱半径修正为 $|x - x_0|$ 或 $|y - y_0|$，如绕 $x = x_0$ 旋转：$\mathrm{d}V = 2\pi |x - x_0| y(x)\,\mathrm{d}x$。

\end{thmbox}
\begin{thmbox}[公式[5]: 已知平行截面面积的立体体积]

若一个立体介于垂直于 $x$ 轴的两平面 $x = a$ 与 $x = b$ ($a < b$) 之间：

\textbf{[1]} 截面积分原理：
过点 $x \in [a, b]$ 作垂直于 $x$ 轴的平行截面，若该截面的面积为已知连续函数 $A(x)$：

取厚度为 $\mathrm{d}x$ 的立体薄片，薄片体积微元为：


\begin{equation*}
\mathrm{d}V = A(x)\,\mathrm{d}x
\end{equation*}


立体的总体积为：


\begin{equation*}
V = \int_a^b A(x)\,\mathrm{d}x
\end{equation*}


\textbf{[2]} 考研典型截面几何形状：
\textbf{(1)}  截面为圆或半圆：$A(x) = \pi R^2(x)$ 或 $\frac{1}{2}\pi R^2(x)$；

\textbf{(2)}  截面为正方形：边长为 $l(x)$，则 $A(x) = l^2(x)$；

\textbf{(3)}  截面为等边三角形：边长为 $l(x)$，则 $A(x) = \frac{\sqrt{3}}{4}l^2(x)$；

\textbf{(4)}  截面为等腰直角三角形：直角边为 $l(x)$，则 $A(x) = \frac{1}{2}l^2(x)$。

\end{thmbox}

\section{3.平面曲线的弧长}

\begin{thmbox}[公式[6]: 三大坐标系下的弧长微元与积分公式]

弧长微元 $\mathrm{d}s$ 是曲线切向位移的欧氏长度度量，满足勾股定理：$(\mathrm{d}s)^2 = (\mathrm{d}x)^2 + (\mathrm{d}y)^2$。

\textbf{[1]} 直角坐标系形式：
\textbf{(1)}  \textit{}显式函数 $y = f(x) \quad (a \leqslant x \leqslant b)$\textit{}：


\begin{equation*}
\mathrm{d}s = \sqrt{(\mathrm{d}x)^2 + (\mathrm{d}y)^2} = \sqrt{1 + [f'(x)]^2}\,\mathrm{d}x
\end{equation*}


弧长为：


\begin{equation*}
s = \int_a^b \sqrt{1 + [f'(x)]^2}\,\mathrm{d}x
\end{equation*}


\textbf{(2)}  \textit{}反函数形式 $x = g(y) \quad (c \leqslant y \leqslant d)$\textit{}：


\begin{equation*}
\mathrm{d}s = \sqrt{1 + [g'(y)]^2}\,\mathrm{d}y \implies s = \int_c^d \sqrt{1 + [g'(y)]^2}\,\mathrm{d}y
\end{equation*}


\textbf{[2]} 参数方程形式：
曲线由 $\begin{cases} x = x(t) \\ \\ y = y(t) \end{cases} \quad (\alpha \leqslant t \leqslant \beta)$ 给出，其中 $x'(t), y'(t)$ 连续且不同时为 0：


\begin{equation*}
\mathrm{d}s = \sqrt{[x'(t)\,\mathrm{d}t]^2 + [y'(t)\,\mathrm{d}t]^2} = \sqrt{[x'(t)]^2 + [y'(t)]^2}\,\mathrm{d}t
\end{equation*}


弧长为：


\begin{equation*}
s = \int_\alpha^\beta \sqrt{[x'(t)]^2 + [y'(t)]^2}\,\mathrm{d}t
\end{equation*}


\textbf{[3]} 极坐标形式：
曲线由 $r = r(\theta) \quad (\alpha \leqslant \theta \leqslant \beta)$ 给出：

由直角坐标与极坐标转换关系：$x = r(\theta)\cos\theta, y = r(\theta)\sin\theta$。

两边微分：


\begin{equation*}
\mathrm{d}x = r'\cos\theta\,\mathrm{d}\theta - r\sin\theta\,\mathrm{d}\theta, \quad \mathrm{d}y = r'\sin\theta\,\mathrm{d}\theta + r\cos\theta\,\mathrm{d}\theta
\end{equation*}


平方求和：


\begin{equation*}
(\mathrm{d}x)^2 + (\mathrm{d}y)^2 = [r^2(\theta) + (r'(\theta))^2]\,(\mathrm{d}\theta)^2
\end{equation*}


故极坐标下的弧长微元为：


\begin{equation*}
\mathrm{d}s = \sqrt{r^2(\theta) + [r'(\theta)]^2}\,\mathrm{d}\theta
\end{equation*}


弧长为：


\begin{equation*}
s = \int_\alpha^\beta \sqrt{r^2(\theta) + [r'(\theta)]^2}\,\mathrm{d}\theta
\end{equation*}


\end{thmbox}

\section{4.旋转曲面的侧面积}

\begin{thmbox}[公式[7]: 旋转曲面侧面积微元与计算公式]

平面光滑弧段绕定轴旋转所生成的曲面称为旋转曲面。

\textbf{[1]} 微元原理与几何剖析：
在曲线上取微元弧段 $\mathrm{d}s$。该小段绕轴旋转一周展开后为一个高（母线长）为 $\mathrm{d}s$、上下底周长近似为 $2\pi r$ 的圆台侧面：


\begin{equation*}
\mathrm{d}S = 2\pi r\,\mathrm{d}s
\end{equation*}


其中 $r$ 为曲线点到旋转轴的垂直距离。

\textbf{核心红线}：微元是弧长 $\mathrm{d}s$，而不是水平投影 $\mathrm{d}x$。用圆柱面积 $2\pi y\,\mathrm{d}x$ 替代会导致系统性误差，因为斜率 $y' \neq 0$ 时 $\mathrm{d}s = \sqrt{1+y'^2}\,\mathrm{d}x > \mathrm{d}x$。

\textbf{[2]} 直角坐标系下的旋转侧面积：
\textbf{(1)}  \textit{}绕 $x$ 轴旋转\textit{}：

回转半径为 $r = |y(x)|$：


\begin{equation*}
S_x = 2\pi \int_a^b |y(x)|\,\mathrm{d}s = 2\pi \int_a^b |y(x)|\sqrt{1 + [y'(x)]^2}\,\mathrm{d}x
\end{equation*}


\textbf{(2)}  \textit{}绕 $y$ 轴旋转\textit{}：

回转半径为 $r = |x|$：


\begin{equation*}
S_y = 2\pi \int_a^b |x|\,\mathrm{d}s = 2\pi \int_a^b |x|\sqrt{1 + [y'(x)]^2}\,\mathrm{d}x
\end{equation*}


\textbf{[3]} 参数方程下的旋转侧面积：
曲线由 $\begin{cases} x = x(t) \\ \\ y = y(t) \end{cases} \quad (\alpha \leqslant t \leqslant \beta)$ 给出：

\textbf{(1)}  \textit{}绕 $x$ 轴\textit{}：


\begin{equation*}
S_x = 2\pi \int_\alpha^\beta |y(t)|\sqrt{[x'(t)]^2 + [y'(t)]^2}\,\mathrm{d}t
\end{equation*}


\textbf{(2)}  \textit{}绕 $y$ 轴\textit{}：


\begin{equation*}
S_y = 2\pi \int_\alpha^\beta |x(t)|\sqrt{[x'(t)]^2 + [y'(t)]^2}\,\mathrm{d}t
\end{equation*}


\textbf{[4]} 极坐标下的旋转侧面积：
曲线 $r = r(\theta) \quad (\alpha \leqslant \theta \leqslant \beta)$ 绕极轴（$x$ 轴，此时纵坐标 $y = r\sin\theta \geqslant 0$）旋转：


\begin{equation*}
S_x = 2\pi \int_\alpha^\beta r(\theta)\sin\theta \sqrt{r^2(\theta) + [r'(\theta)]^2}\,\mathrm{d}\theta
\end{equation*}


\end{thmbox}

\section{5.函数平均值与平面图形形心坐标}

\begin{thmbox}[公式[8]: 连续函数在区间上的平均值]

\textbf{[1]} 定义与几何意义：
设 $f(x)$ 在闭区间 $[a, b]$ 上连续，称数值：


\begin{equation*}
\overline{f} = \frac{1}{b-a}\int_a^b f(x)\,\mathrm{d}x
\end{equation*}


为函数 $f(x)$ 在区间 $[a, b]$ 上的\textbf{算术平均值}（平均值）。

几何意义：由积分第一中值定理，在 $(a, b)$ 内至少存在一点 $\xi$，使得 $\int_a^b f(x)\,\mathrm{d}x = f(\xi)(b-a)$，即 $f(\xi) = \overline{f}$。以区间长 $b-a$ 为底、以平均值 $\overline{f}$ 为高的矩形面积，恰好等于曲边梯形的真实面积。

\textbf{[2]} 离散平均到连续积分的极限过渡：
将区间 $[a, b]$ 进行 $n$ 等分，分点为 $x_i = a + i\frac{b-a}{n}$。采样点函数值的算术平均为：


\begin{equation*}
\frac{1}{n}\sum_{i=1}^n f(x_i) = \frac{1}{b-a}\sum_{i=1}^n f(x_i)\frac{b-a}{n} = \frac{1}{b-a}\sum_{i=1}^n f(x_i)\Delta x
\end{equation*}


当 $n \to \infty$ 时，黎曼和极限即给出连续平均值 $\frac{1}{b-a}\int_a^b f(x)\,\mathrm{d}x$。

\end{thmbox}
\begin{thmbox}[公式[9]: 平面图形与曲线的形心坐标]

形心（几何中心）即假定面密度或线密度均匀恒定（$\rho = 1$）时的质心坐标。

\textbf{[1]} 平面均质薄板的形心 $(\overline{x}, \overline{y})$：
设平面均质区域 $D$ 由 $0 \leqslant y \leqslant f(x), a \leqslant x \leqslant b$ 确定：

\textbf{(1)}  \textit{}总面积 $A$\textit{}：


\begin{equation*}
A = \int_a^b f(x)\,\mathrm{d}x
\end{equation*}


\textbf{(2)}  \textit{}对 $y$ 轴的静矩 $M_y$（决定横坐标 $\overline{x}$）\textit{}：

取宽为 $\mathrm{d}x$ 的竖条，其横坐标为 $x$，微元面积为 $f(x)\,\mathrm{d}x$，静矩微元为 $\mathrm{d}M_y = x\,\mathrm{d}A = x f(x)\,\mathrm{d}x$：


\begin{equation*}
\overline{x} = \frac{M_y}{A} = \frac{\int_a^b x f(x)\,\mathrm{d}x}{\int_a^b f(x)\,\mathrm{d}x}
\end{equation*}


\textbf{(3)}  \textit{}对 $x$ 轴的静矩 $M_x$（决定纵坐标 $\overline{y}$）\textit{}：

竖条的高度为 $f(x)$，由于竖条质量均匀，其形心高度处于中点 $\frac{f(x)}{2}$，静矩微元为 $\mathrm{d}M_x = \frac{f(x)}{2}\,\mathrm{d}A = \frac{1}{2}f^2(x)\,\mathrm{d}x$：


\begin{equation*}
\overline{y} = \frac{M_x}{A} = \frac{\frac{1}{2}\int_a^b f^2(x)\,\mathrm{d}x}{\int_a^b f(x)\,\mathrm{d}x}
\end{equation*}


\textbf{(4)}  \textbf{两曲线夹层区域通式}：

若区域界于 $f_1(x) \leqslant y \leqslant f_2(x)$ 之间：


\begin{equation*}
\overline{x} = \frac{\int_a^b x [f_2(x) - f_1(x)]\,\mathrm{d}x}{\int_a^b [f_2(x) - f_1(x)]\,\mathrm{d}x}, \quad \overline{y} = \frac{\frac{1}{2}\int_a^b [f_2^2(x) - f_1^2(x)]\,\mathrm{d}x}{\int_a^b [f_2(x) - f_1(x)]\,\mathrm{d}x}
\end{equation*}


\textbf{[2]} 平面均质细曲线的质心 $(\overline{x}, \overline{y})$：
设平面连续曲线段为 $L$，总弧长为 $s = \int_L \mathrm{d}s$：


\begin{equation*}
\overline{x} = \frac{\int_L x\,\mathrm{d}s}{\int_L \mathrm{d}s} = \frac{1}{s}\int_L x\,\mathrm{d}s, \quad \overline{y} = \frac{\int_L y\,\mathrm{d}s}{\int_L \mathrm{d}s} = \frac{1}{s}\int_L y\,\mathrm{d}s
\end{equation*}


计算时将 $\mathrm{d}s$ 根据直角坐标、参数方程或极坐标代入对应的微分形式。
\end{thmbox}

\end{document}
